\newpage{}
\section{Creating Sets}

\subsection{Axiom of Existance}
We talked about sets and their equality, but we did not officially say 
there exists a set. However, if we will work with sets, it is crucial.

\begin{center}
    \textit{There exists a set.}
\end{center}
It is a good thing to know that there exists a set, but we do not know 
anything else. Are there multiple sets or there is a unique one, what are 
its elements, and can we sure if there is a set, then all subsets of it are 
also exists?

\subsection{Axiom of Spesification}
If we want to talk about more sets or create new ones using the existing 
sets, we need some other axioms. 
\begin{axiom*}[Axiom of Spesification]
    Suppose A is a set and S(x) is a sentence. Then, there exists a set 
    whose elements are in A for which makes S(x) true.
\end{axiom*}
Informally speaking, this axiom allows us to create subsets for any 
existing set. Now, we can prove some elementary facts about set theory.

\begin{theorem*}[Empty Set] There exists a set with no elements.\end{theorem*}
\begin{proof}
    Since we know there exists a set, let $A$ be a set. Let $S(x)$ be the
    sentence ``$x \not = x$". If we apply the \textit{Axiom of Spesification}
    for $A$ and $S(x)$, the new set must have no elements since
    no $x$ satisfies $x \not= x$.
\end{proof}

\begin{theorem*} There exists a \textbf{unique} set with no elements. \end{theorem*}
\begin{proof}
    Since we know there exists a set with no elements, let $\emptyset_1$ and 
    $\emptyset_2$ be the sets with no elements. According to the 
    \textit{Axiom of Extension}, if these sets are not equal then there 
    must be some element in one of them which the other one does not have.
    But we know that these two sets have no elements. Therefore, they 
    must be equal.
\end{proof}
From now on we will call this set as \textit{empty set} and denote
with \[\emptyset.\]
If $A$ is a set and $S(x)$ is a sentence, we know that we can form a set 
using the \textit{Axiom of Spesification}. Let $B$ be that set.
From now on, we will use the notation \[B = \{x \in A : S(x)\}\]
to say $B$ is a set whose elements are in $A$ for which make $S(x)$ true.

\begin{theorem*} Empty set is a subset of every set.\end{theorem*}
\begin{proof}
  Let $A$ be some set. Suppose that $\emptyset$ is not a subset of $A$.
  Then there exist some $x \in \emptyset $ for which $x \not\in A$.
  Since $\emptyset$ has no elements, this is impossible. Thus, 
  $\emptyset \subseteq A$. Since $A$ was an arbitrary set, we can conclude
  that $\emptyset$ is a subset of every set.
\end{proof}

\begin{theorem*} There are no such set which have all sets as an element.\end{theorem*}
\begin{proof}
    Let $A$ be a set. Using the \textit{Axiom of Spesification} we can
    construct the set $B = \{x \in A : x \notin x\}$. Since $B \in B$ lead 
    us to $B \notin B$, which is a contradiction, $B \notin B$.
    Then, either $B \notin A$ or $B \in B$ (negation of $x \notin x$ for B). 
    Since the latter is impossible, $B \notin A$. $A$ was an arbitrary set.
    Therefore, for all sets, there exists a set which does not belong 
    to the set. Hence, no set can contain all sets.
\end{proof}

This so-called set (a set that contains everything) was called as
\textit{universal set} in early practices of set theory.
This type of sets could cause a paradox named \textit{Russell's Paradox}.
For more information: \url{https://en.wikipedia.org/wiki/Russell's_paradox}


\subsection{Axiom of Pairing}
What we know so far is there exists some set and there is an empty one.
If we want to prove more theorems about sets, we need more axioms.
For example, if we know that $a_0,a_1,\dots,a_n$ are sets, we want
to construct the set $A = \{a_0,a_1,\dots,a_n\}.$ We need some axioms
to achieve this. For start, we will order an axiom to make sure that
for every set $A$ and $B$, there exists a set which contain both $A$ 
and $B$.

\begin{axiom*}[Axiom of Pairing] For all sets A and B, there 
exists a set which contain both $A$ and $B$.\end{axiom*}
Note that this axiom is not saying this set only contains $A$ and $B$,
it may have some other unwanted elements. If we apply 
\textit{Axiom of Spesification} with the sentence $S(x): ``x=A$ or $x=B$",
we get wanted result. \textit{Axiom of Extension} guaranties its 
uniqueness.

If $x$ and $y$ are different sets we will write $\{x,y\}$ to imply 
this is a set which contains only $x$ and $y$ and say this set is 
\textit{unordered pair} or just \textit{pair} of $x$ and $y$.

Since the axiom does not imply that, we can choose our two sets to be equal.
In that case, the set which was created by \textit{Axiom of Pairing} have 
a unique element. Then, we write $\{x\}$ to imply pair of $x$ and $x$ and
say this set is a \textit{singleton} of $x$.

\subsection{Unions and Intersections}
Suppose that $A$ and $B$ are sets. Naturally, we want to unite the 
elements of $A$ and $B$ into a single set.
Since we cannot construct a set based on the axioms we have,
we need some other axioms.

\begin{axiom*}[Axiom of Unions] Suppose that $\mathcal{C}$ is a collection of sets.
\footnote{Actually, everything in the set theory are sets, but if we 
want to emphasise that we will use some properties of sets in some set,
we say \textit{collection of sets} instead ``set of sets". Thus,
this is purely for readability, not because collections and sets are different 
type of objects.} Then, there
exists a set which contains all $x$ for which, $x$ is in at least one element 
of $\mathcal{C}$.\end{axiom*}

Since \textit{Axiom of Unions} says \textit{``which contains"} instead of 
\textit{``only contains"}, sets that are created with \textit{Axiom of Unions}
may have unwanted elements (there are no sets in $\mathcal{C}$ which contains them).
If we use the \textit{Axiom of Spesification}, we can specify the elements 
that are in at least one set in $\mathcal{C}$.

Let $\mathcal{C}$ be a collection of sets. According to the 
\textit{Axiom of Unions} there exists a set $\mathcal{U}$ which have 
all $x$ in at least one set in $\mathcal{C}$. 
Then, $\cup \mathcal{C} = \{x \in \mathcal{U} : 
\text{for some A in } \mathcal{C} \text{ , } x \in A\}$
will be the set whose elements are in some set in $\mathcal{C}$. Similarly,
we can say that if $x \in A$ for some $A \in \mathcal{C}$, then $x \in \cup \mathcal{C}$.
This set $\cup \mathcal{C}$ is called the union of the collection $\mathcal{C}$
of sets, and the \textit{Axiom of Extension} guaranties its uniqueness.

% Will reformat here
Instead of $\cup \mathcal{C}$, some mathematicians use 
\begin{align*}
 \cup \{ X : X \in \mathcal{C} \}   
\end{align*}
or
\begin{align*}
  \cup_{X\in\mathcal{C}} X.
\end{align*}

\begin{theorem*}$\cup{\emptyset} = \emptyset$.\end{theorem*}
\begin{proof}
    Suppose that $x \in \cup \emptyset$. Then, for some $A \in \emptyset$,
    $x \in A$. But, it is impossible to $A \in \emptyset$. Thus, 
    $\cup{\emptyset} = \emptyset$.
\end{proof}

\begin{theorem*}$\cup \{A\} = A$.\end{theorem*}
\begin{proof}
    $(\subseteq)$ Let $x \in \cup\{A\}$. Then, for some $X \in \{A\}$,
    $x \in X$. Since $\{A\}$ have a unique element namely $A$,
    $x \in A$. Thus, $\cup\{A\} \subseteq A$.
   
    $(\supseteq)$ Let $x \in A$. Since $A \in \{A\}$, $x \in \cup\{A\}$.
    Thus, $A \subseteq \cup\{A\}$. 
    
    Since $\cup\{A\} \subseteq A$ and $A \subseteq \cup\{A\}$, 
    $\cup\{A\} = A$.
\end{proof}

\newpage{}

\begin{definition*}
   $A \cup B := \cup\{A,B\}$. 
\end{definition*}
